\documentclass[10pt,aspectratio=169]{beamer}
% Preámbulo modular para presentaciones Beamer (Modelación Estadística)
% Diseño y paleta institucional del Tecnológico de Monterrey

\documentclass[aspectratio=169,xcolor={dvipsnames,table}]{beamer}

\usepackage[utf8]{inputenc}
\usepackage[spanish,mexico]{babel}
\usepackage{amsmath,amssymb,amsfonts,mathtools}
\usepackage{graphicx}
\usepackage{listings}
\usepackage{booktabs}
\usepackage{tikz}
\usetikzlibrary{matrix,arrows,decorations.pathmorphing,shapes.geometric,positioning}

% Paleta Institucional Tecnológico de Monterrey
\definecolor{TecAzul}{HTML}{0039A6}       % Azul TEC: color primario institucional
\definecolor{TecAzulOscuro}{HTML}{1A2E51} % Azul marino oscuro
\definecolor{TecRojo}{HTML}{EC2661}       % Rojo de acento (paleta miTec)
\definecolor{TecGris}{HTML}{646464}       % Gris neutro para comentarios y subtítulos
\definecolor{TecFondoBloque}{HTML}{F4F6F9}% Fondo suave para bloques

% Configuración de Tema Beamer
\usetheme{Boadilla}
\usecolortheme[named=TecAzulOscuro]{structure}
\setbeamercolor{alerted text}{fg=TecRojo}
\setbeamercolor{palette primary}{bg=TecAzulOscuro,fg=white}
\setbeamercolor{palette secondary}{bg=TecAzul,fg=white}
\setbeamercolor{palette tertiary}{bg=TecAzulOscuro!80!black,fg=white}
\setbeamercolor{title}{fg=TecAzulOscuro,bg=TecFondoBloque}
\setbeamercolor{frametitle}{fg=TecAzulOscuro,bg=TecFondoBloque}

% Bloques personalizados elegantes
\setbeamercolor{block title}{bg=TecAzulOscuro,fg=white}
\setbeamercolor{block body}{bg=TecFondoBloque,fg=black}
\setbeamercolor{block title alerted}{bg=TecRojo,fg=white}
\setbeamercolor{block body alerted}{bg=TecRojo!8,fg=black}
\setbeamercolor{block title example}{bg=TecAzul,fg=white}
\setbeamercolor{block body example}{bg=TecAzul!8,fg=black}

% Redondear bordes y añadir sombra sutil
\setbeamertemplate{blocks}[rounded][shadow=true]
\setbeamertemplate{navigation symbols}{}

% Pie de página limpio con número de diapositiva
\setbeamertemplate{footline}{
  \leavevmode%
  \hbox{%
  \begin{beamercolorbox}[wd=.7\paperwidth,ht=2.25ex,dp=1ex,left]{author in head/foot}%
    \hspace*{2ex}\insertshortauthor~~(\insertshortinstitute) \hfill \insertshorttitle
  \end{beamercolorbox}%
  \begin{beamercolorbox}[wd=.3\paperwidth,ht=2.25ex,dp=1ex,right]{date in head/foot}%
    \insertshortdate{}\hspace*{2em}
    \insertframenumber{} / \inserttotalframenumber\hspace*{2ex}
  \end{beamercolorbox}}%
  \vskip0pt%
}

% Configuración de Listings de Python para visualización pedagógica de código
\lstset{
    language=Python,
    basicstyle=\ttfamily\footnotesize,
    keywordstyle=\color{TecAzulOscuro}\bfseries,
    stringstyle=\color{TecRojo},
    commentstyle=\color{TecGris}\itshape,
    showstringspaces=false,
    breaklines=true,
    frame=lines,
    rulecolor=\color{TecAzulOscuro!30},
    backgroundcolor=\color{TecFondoBloque},
    aboveskip=0.5em,
    belowskip=0.5em,
    literate={á}{{\'a}}1 {é}{{\'e}}1 {í}{{\'i}}1 {ó}{{\'o}}1 {ú}{{\'u}}1
             {Á}{{\'A}}1 {É}{{\'E}}1 {Í}{{\'I}}1 {Ó}{{\'O}}1 {Ú}{{\'U}}1
             {ñ}{{\~n}}1 {Ñ}{{\~N}}1 {¿}{{?`}}1 {¡}{{!`}}1
}

% Metadatos institucionales por defecto
\title[Modelación Estadística]{Modelación Estadística para Ciencia de Datos e Ingeniería}
\author[J. Castillo Colmenares]{Juliho David Castillo Colmenares}
\institute[Tec de Monterrey]{Tecnológico de Monterrey \\ Escuela de Ingeniería y Ciencias}
\date{\today}
% Comandos y macros estadísticas compartidas para las presentaciones Beamer (Metropolis)

% Conjuntos numéricos básicos
\newcommand{\R}{\mathbb{R}}
\newcommand{\N}{\mathbb{N}}
\newcommand{\Q}{\mathbb{Q}}
\newcommand{\Z}{\mathbb{Z}}
\newcommand{\set}[1]{\left\{ #1 \right\}}
\newcommand{\abs}[1]{\left|#1\right|}
\newcommand{\norm}[1]{\left\| #1 \right\|}

% Comandos estadísticos y probabilísticos del curso
\newcommand{\Var}{\operatorname{Var}}
\newcommand{\cov}{\operatorname{Cov}}
\newcommand{\corr}{\rho}
\newcommand{\comb}[2]{\begin{pmatrix} #1 \\ #2 \end{pmatrix}}
\newcommand{\s}{\sigma}
\newcommand{\particion}{\mathfrak{P}}
\newcommand{\card}[1]{\#\left( #1 \right)}
\newcommand{\seq}[3]{\set{#1}_{#2}^{#3}}
\newcommand{\E}{\mathbb{E}}
\newcommand{\Prob}{\mathbb{P}}

% Entornos pedagógicos y cajas en español académico natural
\newenvironment{discusion}{%
  \begin{alertblock}{Pregunta de Discusión}%
}{%
  \end{alertblock}%
}

\newenvironment{reflexion}{%
  \begin{alertblock}{Reflexión para el Aula}%
}{%
  \end{alertblock}%
}

\newenvironment{paradoja}{%
  \begin{alertblock}{Paradoja de Exploración}%
}{%
  \end{alertblock}%
}

\newenvironment{motivacion}{%
  \begin{exampleblock}{Motivación: Ciencia de Datos en la Práctica}%
}{%
  \end{exampleblock}%
}

\newenvironment{retoclase}{%
  \begin{block}{Reto Rápido en Equipo}%
}{%
  \end{block}%
}

\title[Función generadora de momentos]{Sección 04.06: Función generadora de momentos}\subtitle{Momentos, unicidad y sumas independientes}\author[J. Castillo Colmenares]{Juliho Castillo Colmenares}\institute{Tecnológico de Monterrey}\date{}
\begin{document}
\begin{frame}[plain]\titlepage\end{frame}
\begin{frame}{Hoja de ruta}\small\begin{enumerate}\item Definición y existencia.\item Momentos, unicidad y suma.\item Ejemplos y tabla de referencia.\end{enumerate}\end{frame}
\begin{frame}{Fuente y alcance}\small La función generadora de momentos (FGM) codifica una distribución en una sola función y conecta cálculos de momentos con propiedades de sumas independientes.\begin{block}{Pregunta guía}¿Qué información de una distribución queda almacenada en $M_X(t)$?\end{block}\end{frame}
\begin{frame}{Definición}\small\[M_X(t)=E[e^{tX}].\]La FGM transforma la variable aleatoria en una función del parámetro $t$.\end{frame}
\begin{frame}{Caso discreto y continuo}\small\begin{itemize}\item Discreto: $M_X(t)=\sum_x e^{tx}f(x)$.\item Continuo: $M_X(t)=\int_{-\infty}^{\infty}e^{tx}f(x)\,dx$.\end{itemize}\end{frame}
\begin{frame}{Existencia}\small La FGM existe cuando la suma o integral converge en un intervalo abierto alrededor de $t=0$. No toda distribución tiene FGM en un entorno de cero.\end{frame}
\begin{frame}{Momentos mediante derivadas}\small Si existe la FGM,\[E[X^n]=M_X^{(n)}(0)=\left.\frac{d^nM_X(t)}{dt^n}\right|_{t=0}.\]\end{frame}
\begin{frame}{Primeros momentos}\small\[E[X]=M_X'(0),\qquad E[X^2]=M_X''(0),\qquad \operatorname{Var}(X)=M_X''(0)-[M_X'(0)]^2.\]\end{frame}
\begin{frame}{Idea de la demostración}\small Bajo condiciones de regularidad, se intercambian derivación y esperanza:\[M_X^{(n)}(t)=E[X^ne^{tX}].\]Evaluar en $t=0$ elimina el factor exponencial.\end{frame}
\begin{frame}{Propiedad de unicidad}\small Si dos variables tienen la misma FGM en un intervalo alrededor de cero, tienen la misma distribución. La FGM puede identificar la ley cuando existe.\end{frame}
\begin{frame}{Propiedad de suma}\small Para variables independientes $X$ y $Y$,\[M_{X+Y}(t)=M_X(t)M_Y(t).\]La independencia convierte la FGM de una suma en un producto.\end{frame}
\begin{frame}{Idea de la demostración de suma}\small\[E[e^{t(X+Y)}]=E[e^{tX}e^{tY}]=E[e^{tX}]E[e^{tY}].\]El segundo paso usa independencia.\end{frame}
\begin{frame}{Ejemplo Bernoulli: planteamiento}\small Si $X$ es Bernoulli con probabilidad de éxito $p$, solo toma $0$ y $1$. Por tanto,\[M_X(t)=(1-p)+pe^t.\]\end{frame}
\begin{frame}{Ejemplo Bernoulli: momentos}\small\[M_X'(0)=p,\qquad M_X''(0)=p,\qquad\operatorname{Var}(X)=p(1-p).\]La FGM recupera los momentos conocidos.\end{frame}
\begin{frame}{Ejemplo normal: planteamiento}\small Para $Z\sim N(0,1)$, completar el cuadrado en la integral produce una exponencial cuadrática.\end{frame}
\begin{frame}{Ejemplo normal: resultado}\small Si $X\sim N(\mu,\sigma^2)$,\[M_X(t)=\exp\!\left(\mu t+\frac{\sigma^2t^2}{2}\right).\]\end{frame}
\begin{frame}{Ejemplo exponencial: planteamiento}\small Si $X\sim\operatorname{Exp}(\lambda)$, la integral converge para $t<\lambda$ y\[M_X(t)=\lambda\int_0^\infty e^{-(\lambda-t)x}\,dx.\]\end{frame}
\begin{frame}{Ejemplo exponencial: resultado}\small\[M_X(t)=\frac{\lambda}{\lambda-t}=\left(1-\frac t\lambda\right)^{-1}.\]Además, $E[X^n]=n!/\lambda^n$.\end{frame}
\begin{frame}{Tabla de referencia}\small\begin{itemize}\item Bernoulli: $(1-p)+pe^t$.\item Binomial: $(1-p+pe^t)^n$.\item Poisson: $\exp(\lambda(e^t-1))$.\item Normal: $\exp(\mu t+\sigma^2t^2/2)$.\end{itemize}\end{frame}
\begin{frame}{Aplicación de la suma}\small Si $X_1,\ldots,X_n$ son ensayos Bernoulli independientes con probabilidad $p$,\[M_{\sum X_i}(t)=(1-p+pe^t)^n,\]que identifica la distribución binomial.\end{frame}
\begin{frame}{Interpretación}\small Las derivadas convierten una función compacta en momentos; la multiplicación convierte independencia en una herramienta para construir distribuciones de sumas.\end{frame}
\begin{frame}{Síntesis y conexión}\small La FGM ofrece definición, cálculo de momentos, unicidad y suma. Estas propiedades enlazan distribuciones discretas, continuas y modelos compuestos.\end{frame}
\end{document}
